% Required packages; load this file before \begin{document}.
\usepackage{amsmath,amssymb,amsthm}
\usepackage{longtable,booktabs,array}
\usepackage{tikz}
\usetikzlibrary{arrows.meta,positioning,shapes.geometric,fit,calc}
\usepackage{hyperref}
\renewcommand{\proofname}{证明}
\newtheorem{modelprop}{模型命题}[chapter]
\newtheorem{modeldef}{模型定义}[chapter]
\newcommand{\lawref}[2]{\textsuperscript{\hyperlink{#1}{[#1：#2]}}}
\newcommand{\lawinput}[1]{\par\noindent\textbf{法律前提。}#1\par}
\newcommand{\modelbound}[1]{\par\noindent\textbf{适用边界。}#1\par}
\newcommand{\memorycard}[1]{\par\noindent\textbf{闭卷检索卡。}#1\par}
\newcommand{\caseinput}[1]{\par\noindent\textbf{教学设例。}#1\par}
\newcommand{\formalstatus}{\par\noindent\textbf{版本与范围：}依据上传书目录的二十五讲建立核心学习模型；这是独立编写的形式化研习稿，不能据此声称原书每一段、每一道例题均已完成法律校订或形式证明。法律核验记录日期为2026年11月20日，具体核验范围见本册末尾。}

\providecommand{\llbracket}{\mathopen{[\![}}
\providecommand{\rrbracket}{\mathclose{]\!]}}
